\chapter{A complex is an operation}
\label{ch:operads}

% Part V, added 2026-09-25 from ~/Downloads/chygraph_master_equation/draft.tex,
% Sec. 8. No numbers are computed here. The genus counts in Sec. sec:op-modular
% are the first Betti number of the incidence structure, arithmetic on objects
% the book has already drawn (Figs. fig:sharededge and fig:ring).

Chapter~\ref{ch:chyequation} wrote every interior sum in the book as one
operator, $\Tint_{a}$, that takes a message on each leg of a complex but one
and returns a message on the leg left open. An object with $c$ inputs and one
output, together with a rule for plugging the output of one into an input of
another, is what algebra calls an \emph{operation} in an \emph{operad}. This
chapter says that the identification is exact rather than suggestive, and
reads three things off it: why percolation collapses to generating functions,
what the excess bracket is called, and why Part~IV is a different kind of
object from Parts~II and~III.

The chapter is short because the correspondence is direct. What it adds to
the earlier chapters is vocabulary rather than results: a name for the
composition the two steps of Chapter~\ref{ch:statmech} perform, a name for
the discount rule of Fig.~\ref{fig:themap}, and a name --- genus --- for the
count of loops over complexes that Chapter~\ref{ch:overlap} found organises
the errors. Chapter~\ref{ch:mobius} gives the same count a homological
reading.

\section{The connection}
\label{sec:op-connection}
\index{operad|textbf}
\index{interior operator}

An operad is a family of operations, one set $\mathcal{O}(n)$ for each arity
$n$, each operation having $n$ inputs and one output, with a composition rule
for feeding outputs into inputs and an action of the symmetric group
permuting the inputs \citep{markl2002}. An \emph{algebra} over an operad is a
vector space $V$ on which every operation of arity $n$ acts as a multilinear
map $V^{\otimes n}\to V$, the maps respecting the compositions. The
definition asks for nothing else, and every piece of it is already in
Chapter~\ref{ch:chyequation}.

The space $V$ is the message space, functions on $\Omega_{i}$ up to a positive
factor, or laws over such functions at the level of the ensemble. The
operation of arity $c-1$ attached to a layer-$l$ complex is its interior
operator: Eq.~\eqref{eq:chyop}'s $\Tint^{(i)}_{a}$ takes the $c-1$ member
messages on the legs other than $i$ and returns a message on leg $i$, and
Eq.~\eqref{eq:act} lets the same operator act on laws by pushforward. The
symmetric-group action is the symmetry of $\Tint_{a}$ under relabelling the
members, which is what made the pushforward in Eq.~\eqref{eq:act} well
defined without an ordering. And composition is nesting: a complex $a$ whose
member $j$ is itself a complex has, on leg $j$, not a bare variable but the
output of $\Tint_{j}$ acting on $j$'s own members, Eq.~\eqref{eq:chyup}.
Plugging one interior operator into an input of another is the down step of
one level feeding the up step of the next, and that is the only composition
the recursion ever performs. The family $(\Tint_{n})_{n\ge0}$ of symmetric
multilinear maps in Eq.~\eqref{eq:act} is therefore an algebra over a
symmetric operad, and the two steps of Sec.~\ref{sec:twosteps} are its
structure maps.

There is a second operation in the recursion, and it is the one with no model
in it. The exterior product $\otimes$ of Sec.~\ref{sec:ce-bp} --- pointwise
multiplication of functions on one variable, addition of cavity fields ---
is commutative and associative and has one operation in every arity: the
product of $n$ things. The operad with exactly that content is $\mathrm{Com}$,
the commutative operad, and the up step Eq.~\eqref{eq:up} is its action.
So the recursion is an algebra over an operad generated by two kinds of
operation: $\mathrm{Com}$ for the relay at a vertex, and one interior
operation per layer for the model inside a complex. The Hamiltonian sits in
the second and nowhere else, which is Sec.~\ref{sec:ce-bp}'s sentence again.

\paragraph{Percolation is $\mathrm{Com}$.} The first row of
Table~\ref{tab:ce-substitutions} said that for percolation the interior
operator is itself a product: reachability factorises, so
$(T_{n})\push\,\delta_{Q}^{\otimes n}=\delta_{Q^{n}}$. In the present language
the interior operation of every layer is an operation of $\mathrm{Com}$, the
same operad the exterior product already lives in. An algebra over
$\mathrm{Com}$ is a commutative algebra, in which composing operations is
multiplying, and multiplying point masses composes the generating functions
that weight them. That is why Eqs.~\eqref{eq:Qminus}--\eqref{eq:Qplus} are
compositions of generating functions and nothing more, and why the whole of
Part~II fits in $2L^{2}$ numbers: there is no operation in the calculation
that is not a product. It is also why Chapter~\ref{ch:potts}'s $q\to1$ limit
lands where it does. Section~\ref{sec:qlimit} showed that as $q\to1$ the
ratio a Potts complex emits becomes a product over its members, and the
interior sum becomes the intra-complex generating function $\bar G$: at $q=1$
the operation the complex performs degenerates to the one $\mathrm{Com}$
already has, and the generating function is what is left of it. Every other
row of Table~\ref{tab:ce-substitutions} is an operation outside
$\mathrm{Com}$, and is the reason the message stopped being a number.

\section{Partial composition and the excess bracket}
\label{sec:op-partial}
\index{excess!generating function}
\index{partial composition|textbf}

Operads compose in two ways that are equivalent but read differently. Full
composition plugs an operation into every input of another at once. Partial
composition, written $p\circ_{i}q$, plugs $q$ into the $i$th input of $p$
alone, leaving the other inputs of $p$ and all of those of $q$ free: an
operation of arity $n$ composed on its $i$th input with one of arity $m$ has
arity $n+m-1$. The $-1$ is the input consumed by the composition.

That $-1$ is the excess bracket. Figure~\ref{fig:themap}'s rule was
``discount the inclusion you came in by, and only from routes back into that
layer''. Read as a composition: the complex arrived at is $p$, the inclusion
arrived along is the input $i$ being composed on, and the free legs of $p$
after composition are its other inputs. Since the composed-on leg belongs to
one definite layer, it is subtracted from that layer's count and from no
other, which is exactly what $[\bar G^{l}_{k}]_{k=m}$ does: the excess
function on the layer of arrival, the ordinary one on the rest. The same
reading serves the up direction. A vertex reached from a container is the
operation $\otimes$ of $\mathrm{Com}$ composed on one of its inputs, and the
containers left free are counted by $[\bar\Phi^{l}_{k}]_{k=m}$. The two
excess families of Sec.~\ref{sec:ensembles}, one per direction, are the two
places a partial composition can consume a leg.

Equation~\eqref{eq:chyop} names the leg differently, and the difference is
worth stating. There the superscript $(i)$ on $\Tint^{(i)}_{a}$ names the leg
left \emph{open}, the output. In an ordinary operad an operation has one
output, fixed once and for all; on a chygraph any member can be the one
answered, and $u_{a\to i}$ makes $i$ the output while $h_{i\to a}$ composes
on $i$ as an input. Both are fixed by one choice, and the choice is the
directed inclusion the message runs along. An operad in which every leg can
serve as the output is a \emph{cyclic} operad, and that is the subject of the
next section.

\section{Cyclic and modular}
\label{sec:op-modular}
\index{cyclic operad|textbf}
\index{modular operad|textbf}
\index{genus|textbf}
\index{overlap!between complexes}

A cyclic operad forgets which leg is the output: an operation of arity $n$
becomes an object with $n+1$ legs, any of which may be composed on, and
compositions are trees whose vertices are operations and whose edges are the
legs joined. Messages running both ways along every inclusion, with any leg
able to serve as the output, is that structure, and the treelike theory of
Parts~II and~III is cyclic: every composition the recursion performs is a
tree, because that is what treelike over complexes, Sec.~\ref{sec:treelike},
means.

\citet{getzler1998} define a \emph{modular} operad by adding one more
composition: a \emph{self-gluing} $\xi_{ij}$ that joins two legs of the same
operation to each other. Compositions are then graphs rather than trees, and
the operations carry a second grading, the \emph{genus} $g$, which a
self-gluing raises by one. The genus of a composition is the first Betti
number of the graph it draws: edges minus vertices plus components.

Part~IV is modular. Two complexes sharing two atoms, Fig.~\ref{fig:sharededge},
are two operations joined along two legs, which is one composition followed
by one self-gluing. On the incidence structure --- four atoms, two complexes,
six inclusions --- the count is $6-6+1=1$: genus one. The ring of $k$
triangles glued at single atoms, Fig.~\ref{fig:ring}, has $2k$ atoms, $k$
complexes and $3k$ inclusions, so $3k-3k+1=1$: genus one again, though no
pair of complexes shares two atoms. The two objects Chapters~\ref{ch:overlap}
and~\ref{ch:metacomplex} used as their smallest counterexamples are the two
smallest self-gluings, one made in a single composition and one made around a
cycle, and the recursion double-counts on both for the same reason: it is a
cyclic calculation run on a modular composition.

\paragraph{What merging does.} The merge closure of Sec.~\ref{sec:merging}
takes the two triangles of Fig.~\ref{fig:sharededge} and replaces them by
one complex of four atoms: one operation of arity four in place of two of
arity three glued twice. On the incidence structure that is four atoms, one
complex, four inclusions, $4-5+1=0$. Merging absorbs a self-gluing into a
single operation, lowering the genus by one, and sums the loop it absorbed
inside the interior, where Chapter~\ref{ch:statmech} pays for it at a cost
set by the cardinality. It cannot touch the ring, because the ring's
self-gluing is not between two operations but around $k$ of them, and no
pairwise rule sees it; that is Sec.~\ref{sec:mergeerror}'s finding, and the
reason Chapter~\ref{ch:exactness} had to go to the incidence structure as a
whole. The exactness condition of that chapter, in this vocabulary, is
\emph{genus zero after merging}: every composition the ensemble performs is a
tree of operations, whatever their arities.

That the overlap corrections of Chapter~\ref{ch:overlap} are organised by
the number of independent loops over complexes is what the modular grading
predicts. Getzler and Kapranov's construction is graded by genus because
that is the quantity a self-gluing changes; a calculation exact at genus zero
and wrong above it will have its error indexed the same way.
Chapter~\ref{ch:mobius} makes the same count as a first Betti number and
attaches a systematic expansion to it.

\begin{figure}[t]
\centering
\begin{adjustbox}{width=0.88\linewidth}
\begin{tikzpicture}[x=1cm,y=1cm]
% ---- left: a tree composition (cyclic) ----
\begin{scope}
  \node[atm] (a1) at (0,0) {};
  \node[atm] (a2) at (0.9,0) {};
  \node[atm] (a3) at (1.8,0) {};
  \node[atm] (a4) at (2.7,0) {};
  \node[atm] (a5) at (3.6,0) {};
  \node[cpx] (A) at (0.9,1.2) {};
  \node[cpx] (Bc) at (2.7,1.2) {};
  \draw[inc] (a1)--(A) (a2)--(A) (a3)--(A);
  \draw[inc] (a3)--(Bc) (a4)--(Bc) (a5)--(Bc);
  \draw[msg] ($(a3)+(-0.12,0.2)$) -- ($(A)+(0.28,-0.18)$);
  \draw[msg] ($(Bc)+(-0.3,-0.18)$) -- ($(a3)+(0.12,0.22)$);
  \node[lb,left=2pt of A] {$p$};
  \node[lb,right=2pt of Bc] {$q$};
  \node[lb,below=1pt of a3] {$i$};
  \node[tb,anchor=north,align=center] at (1.8,-0.55)
    {cyclic: $p\circ_{i}q$, a tree\\genus $0$};
\end{scope}
% ---- right: a self-gluing (modular) ----
\begin{scope}[xshift=6.2cm]
  \node[atm] (b1) at (0,0) {};
  \node[atm] (b2) at (0.9,0) {};
  \node[atm] (b3) at (1.8,0) {};
  \node[atm] (b4) at (2.7,0) {};
  \node[cpx] (C) at (0.9,1.2) {};
  \node[cpx] (D) at (1.8,1.2) {};
  \draw[inc] (b1)--(C) (b2)--(C) (b3)--(C);
  \draw[inc] (b2)--(D) (b3)--(D) (b4)--(D);
  \draw[incd,line width=0.8pt] (b2) -- (C) -- (b3) -- (D) -- (b2);
  \node[lb,left=2pt of C] {$p$};
  \node[lb,right=2pt of D] {$q$};
  \node[tb,anchor=north,align=center] at (1.35,-0.55)
    {modular: glued on two legs\\genus $1$};
\end{scope}
\end{tikzpicture}
\end{adjustbox}
\caption{\textbf{Two ways to compose.} Left: two complexes meeting at one
atom $i$, drawn as the inclusion diagram of Sec.~\ref{sec:drawing}. As
operations this is the partial composition $p\circ_{i}q$: the leg $i$ is
consumed, the two messages on it run in opposite directions, and the
composition is a tree. Every composition in Parts~II and~III has this form.
Right: two complexes sharing two atoms, Fig.~\ref{fig:sharededge} again. A
composition followed by a self-gluing, in the sense of \citet{getzler1998}:
the dashed cycle complex--atom--complex--atom is the loop the gluing creates,
its genus is one, and the cyclic recursion double-counts it. The merge closure
of Chapter~\ref{ch:metacomplex} replaces the right panel by one operation of
arity four and returns the genus to zero.}
\label{fig:opgluing}
\end{figure}

\section{Compound complexes are properads}
\label{sec:op-properads}
\index{properad|textbf}
\index{compound complex}
\index{intra-complex component}

One thing in Chapter~\ref{ch:chygraphs} does not fit in an operad, and it is
worth saying where it fits instead. Section~\ref{sec:ensembles} distinguished
a complex's cardinality $c$ from the size $s$ of the intra-complex component
entered, because the interior hypergraph of a complex may be several
disjoint pieces: a publication entered as one of its authors reaches the
other authors and not the references. Chapter~\ref{ch:giant}'s constructions,
Sec.~\ref{sec:constructions}, keep the two apart throughout, and the
threshold tensor carries $\ave{s}$ and $\ave{\sbar}$ rather than the
cardinality.

An operation has one output. A compound complex entered through one
component emits into that component, and entered through another emits into
that one; it is, as far as the members are concerned, several operations
that happen to share a container and a state. Once the container messages
and the sum over $\sigma_{a}$ of Eq.~\eqref{eq:chydown} are included the
pieces are coupled, so the complex is not simply a disjoint union of
operations either. The object with several inputs and several outputs,
composed along several legs at once, is a \emph{properad}, an operad's
generalisation in which operations have as many outputs as inputs they
please, and the natural home for a compound complex is there: one output per
interior component, coupled through the complex's own leg. This is offered as
the right vocabulary and not as a theorem; nothing in the book depends on it,
because Part~II wrote the percolation theory of compound chygraphs directly
in terms of $s$ and never needed the name. What the name says is that the
percolation theory of random properads is already done, in
Chapters~\ref{ch:percolation} and~\ref{ch:giant}, for the class whose
operations are $\mathrm{Com}$.

\section{What it gives the book}
\label{sec:op-gives}

A name for the two halves of the book's difficulty. The treelike theory of
Parts~II and~III is cyclic: every complex has one output at a time and the
compositions are trees, which is what makes the recursion close on one law
per directed inclusion type. Part~IV, where two complexes share two atoms and
the recursion double-counts, is modular: the composition has a self-gluing,
and the grading that measures it is the genus, the first Betti number of the
incidence structure. Chapter~\ref{ch:exactness}'s condition is genus zero
after merging.

A name for the discount rule. The excess bracket, explained in
Chapter~\ref{ch:percolation} by a figure and a paragraph, is partial
composition: the leg composed on is consumed, and it is consumed from one
layer's count because it belongs to one layer. Chapter~\ref{ch:species} gives
the same bracket a second name, a derivative, and the two agree because a
partial composition on the $i$th input and a derivative with respect to the
sort of that input both remove one slot.

And a reason, rather than an observation, for the fact that percolation is
the case that compresses. The interior operation of percolation is an
operation of $\mathrm{Com}$, the operad the exterior product already lives
in, so the calculation contains nothing but products and every composition
of products is a composition of generating functions. Any interaction that
is not a product leaves $\mathrm{Com}$, and the message stops being a number
at that moment and for that reason.

\section{What the field could get}
\label{sec:op-field}
\index{random operad|textbf}

Operads meet probability rarely, and this is where the chygraph offers most.

The ensemble of Sec.~\ref{sec:ensembles} is a probability measure on the
compositions of a free cyclic operad generated by one operation per layer,
with the chy-degree and intra-complex generating functions fixing how many
legs each operation has and which operations it is composed with. The
chygraph equation, Eqs.~\eqref{eq:ensplus}--\eqref{eq:ensminus}, is then the
characteristic equation of that measure: the law of the message on a random
leg reproduces itself under composition. That gives operadic composition a
statistical mechanics it does not have. A free energy per operation is
Eq.~\eqref{eq:bethe}, whose counting numbers are one per operation, one per
leg and minus one per composition. A threshold is the Perron root of the
linearised composition crossing one, Eq.~\eqref{eq:branch}, and a critical
amplitude is its second derivative, Sec.~\ref{sec:amplitude}. All of them are
read off the four moment matrices of Eq.~\eqref{eq:moments}, which are the
first moments of the arities and of the composition counts.

Two structural points follow. The excess bracket is partial composition, so
the discount rule of the physics has an operadic name and the operadic
composition has a probabilistic one. And a compound complex is an operation
with several outputs, so compound chygraphs are random properads, and
Part~II already contains their percolation theory.

On the modular side, the overlap density controls the genus, and the
question of when random compositions acquire loops is a percolation
transition in the operad. Chapter~\ref{ch:exactness} located it on the
ensembles the book uses: never reached for placed motifs, at any branching
ratio; never reached in one dimension, at any density; reached at a
definite density in two and three. In this vocabulary
that chapter is the statement that a random composition in a free cyclic
operad stays a tree below a threshold and acquires extensive genus above it,
with the threshold fixed by the same four matrices. Whether the modular
operad's own structure --- the grading by genus, the self-gluing as a
composition in its own right --- gives a way to organise the calculation
above the threshold, which the Outlook lists as the largest thing the book
does not do, is not known. It is the natural place to look.
